\section{Fourier coordinates for regional parent spaces}
\label{sec:peps_fourier_parent_transport}

Let $G$ be finite, and let $U,L$ be representations on finite coordinate
sets $X,Y$. Suppose $Q:\C^X\to\C^Y$ satisfies
$Q^\dagger Q=I$ and $L(g)=QU(g)Q^\dagger$.
All graph edges retain their ordered tail--head orientation.
The physical bond coordinate map is $B(Q)=Q\otimes\overline Q$ in
head--tail order, as in \cite[Section~7]{Schuch2010PEPS}.
Every open contraction retains its full virtual boundary coefficient.

\subsection{Regional coordinate changes}

\begin{theorem}[The reverse coordinate identity]%
    \label{thm:peps_fourier_reverse_coordinates}
    \lean{TNLean.PEPS.mul_conjTranspose_eq_one_of_representation_coordinates,
        TNLean.PEPS.representation_coordinates_conjTranspose}
    \leanok
    The identity $L(g)=QU(g)Q^\dagger$ alone implies $QQ^\dagger=I$.
    If also $Q^\dagger Q=I$, then $U(g)=Q^\dagger L(g)Q$.
\end{theorem}
\begin{proof}\leanok
    Evaluate the first identity at $g=1$, since $U(1)=L(1)=I$.
    For the second, multiply it on the left by $Q^\dagger$ and on
    the right by $Q$, and cancel $Q^\dagger Q$.
\end{proof}

\begin{definition}[Oriented regional Fourier matrices]%
    \label{def:peps_fourier_regional_matrix}
    \lean{TNLean.PEPS.graphBoundaryCoordinateMatrix,
        TNLean.PEPS.graphRegionCoordinateMatrix}
    \leanok
    For a crossing edge $f=(t,h)$ of a region $R$, put
    $Q_f=\overline Q$ when $t\in R$ and $Q_f=Q$ when $h\in R$.
    On its physical crossing endpoints and internal head--tail pairs,
    define
    \begin{align}
        Q_R=\bigotimes_{f\in\partial E_R}Q_f
        \otimes\bigotimes_{f\in E_R}B(Q).
        \label{eq:peps_fourier_regional_matrix}
    \end{align}
    The virtual boundary matrix is
    $Q_{\partial R}=\bigotimes_{f\in\partial E_R}Q_f$.
\end{definition}

\begin{theorem}[Composition of product coordinate matrices]%
    \label{thm:peps_fourier_product_composition}
    \lean{TNLean.PEPS.mixedPhysicalProductMatrix_mul,
        TNLean.PEPS.mixedPhysicalProductMatrix_one,
        TNLean.PEPS.bondCoordinateMatrix_mul,
        TNLean.PEPS.bondCoordinateMatrix_one,
        TNLean.PEPS.graphRegionCoordinateMatrix_mul_conjTranspose}
    \leanok
    \uses{def:peps_fourier_regional_matrix}
    For finite products with compatible index spaces, multiplying
    product matrices multiplies the matrices at each factor, and a
    product of identity matrices is the identity. In particular,
    $B(P)B(Q)=B(PQ)$ and $B(I)=I$.
    If $QQ^\dagger=I$, then $Q_R(Q^\dagger)_R=I$.
\end{theorem}
\begin{proof}\leanok
    In matrix multiplication, the sum over an intermediate product
    configuration factors into the product of the individual sums.
    Conjugation commutes with multiplication, proving the bond formula.
    For a tail crossing the two factors are $\overline Q$ and
    $\overline{Q^\dagger}$, with product
    $\overline{QQ^\dagger}=I$; at a head they are $Q,Q^\dagger$.
    The internal product is $B(QQ^\dagger)=I$, proving the regional
    cancellation identity.
\end{proof}

\begin{theorem}[Fourier transformation of open generators]%
    \label{thm:peps_fourier_open_contractions}
    \lean{TNLean.PEPS.graphOpenBoundaryFactor_coordinates,
        TNLean.PEPS.graphOpenInternalFactor_coordinates,
        TNLean.PEPS.graphRegionCoordinateMatrix_open,
        TNLean.PEPS.graphRegionCoordinateMatrix_open_superposition}
    \leanok
    \uses{def:peps_fourier_regional_matrix}
    Let $a_R^U(\theta)$ and $a_R^L(\eta)$ be the actual unweighted
    open averaging contractions in regional bond coordinates.
    If $Q^\dagger Q=I$ and $L(g)=QU(g)Q^\dagger$, then
    \begin{align}
        Q_Ra_R^U(\theta)
                            & =\sum_\eta\left(\prod_{f\in\partial E_R}
        Q_f(\eta_f,\theta_f)\right)a_R^L(\eta),
        \label{eq:peps_fourier_open_generator}                         \\
        Q_R\mathcal C_R(U)c & =\mathcal C_R(L)Q_{\partial R}c.
        \label{eq:peps_fourier_open_superposition}
    \end{align}
    Here $c$ is an arbitrary virtual boundary coefficient vector.
    For each internal edge the individual factor is transformed by
    $B(Q)$, using only $L(g)=QU(g)Q^\dagger$.
    At each crossing the individual factor obeys the corresponding
    one-ended sum over $\eta_f$ in
    (\ref{eq:peps_fourier_open_generator}).
\end{theorem}
\begin{proof}\leanok
    For an internal edge, vectorization of
    $L(q_hq_t^{-1})=QU(q_hq_t^{-1})Q^\dagger$ gives its factor identity.
    At a head crossing use $L(g)Q=QU(g)$; at a tail crossing use
    $Q^\dagger L(g)=U(g)Q^\dagger$, with $g=q_t^{-1}$, and take
    coefficients in tail order. These identities leave the oriented
    virtual boundary matrix $Q_f$.
    Expand the site averages and distribute the product of crossing
    sums to obtain~(\ref{eq:peps_fourier_open_generator}).
    Summing against $c(\theta)$ and interchanging finite sums proves
    (\ref{eq:peps_fourier_open_superposition}).
\end{proof}

\begin{theorem}[Exact equality of every regional range]%
    \label{thm:peps_fourier_regional_ranges}
    \lean{TNLean.PEPS.graphRegionCoordinateMatrix_maps_openBondSpace,
        TNLean.PEPS.graphRegionCoordinateMatrix_map_openBondSpace}
    \leanok
    Under the coordinate hypotheses above,
    $Q_R\mathcal G_R(U)=\mathcal G_R(L)$ in regional bond coordinates,
    for every $R$. In particular the forward range inclusion holds.
\end{theorem}
\begin{proof}\leanok
    \uses{thm:peps_fourier_open_contractions,
        thm:peps_fourier_reverse_coordinates,
        thm:peps_fourier_product_composition}
    Equation~(\ref{eq:peps_fourier_open_generator}) maps every source
    generator into the target span. Apply it also to $Q^\dagger$,
    using Theorem~\ref{thm:peps_fourier_reverse_coordinates}, to send
    every target range vector back into the source range.
    The identity $Q_R(Q^\dagger)_R=I$ then supplies its preimage.
\end{proof}

\subsection{Exterior blocks and complete kernels}

\begin{definition}[Exterior crossing coefficients]%
    \label{def:peps_fourier_exterior_coefficient}
    \lean{TNLean.PEPS.graphBoundaryExteriorCoordinateCoefficient}
    \leanok
    For fixed exterior crossing configurations $\gamma_X,\gamma_Y$,
    define
    \begin{align}
        c_{\partial R}(\gamma_Y,\gamma_X)
        =\prod_{f=(t,h)\in\partial E_R}
        \begin{cases}
            Q(\gamma_{Y,f},\gamma_{X,f}),            & t\in R, \\
            \overline{Q(\gamma_{Y,f},\gamma_{X,f})}, & h\in R.
        \end{cases}
        \label{eq:peps_fourier_exterior_coefficient}
    \end{align}
    The exterior endpoint has the complementary orientation to the
    retained regional endpoint.
\end{definition}

\begin{theorem}[The actual global Fourier blocks preserve local ranges]%
    \label{thm:peps_fourier_exterior_blocks}
    \lean{TNLean.PEPS.mixedPhysicalProductMatrix_boundary_coordinates,
        TNLean.PEPS.regionOperatorBlock_graphNumberedBondMatrix_coordinates_maps,
        TNLean.PEPS.graphNumberedBondMatrix_coordinates_maps_regionParentGroundSpace}
    \leanok
    \uses{def:peps_fourier_exterior_coefficient}
    Let $\mathcal Q$ be $\bigotimes_{f\in E}B(Q)$ in numbered site
    coordinates. For fixed input and output exterior configurations,
    its regional matrix block is the product of the scalar from
    wholly exterior bonds, the coefficient
    (\ref{eq:peps_fourier_exterior_coefficient}), and $Q_R$ in numbered
    regional coordinates. The same factorization holds for the
    mixed product of its crossing and internal matrices before numbering.
    Hence every block maps $\mathcal G_R(U)$ into $\mathcal G_R(L)$,
    and $\mathcal Q\mathcal P_{\mathcal R}(U)
        \subseteq\mathcal P_{\mathcal R}(L)$ for any region family.
\end{theorem}
\begin{proof}\leanok
    \uses{thm:peps_fourier_regional_ranges,thm:peps_parent_support_bond_blocks,
        thm:peps_parent_support_open_coordinates,thm:peps_parent_support_block_criterion}
    Fixing one endpoint in
    $B(Q)_{(y_h,y_t),(x_h,x_t)}=Q_{y_h,x_h}\overline{Q_{y_t,x_t}}$
    extracts exactly the exterior scalar and leaves $Q_f$ on the
    regional endpoint. Multiplying over edges gives the block formula.
    Theorem~\ref{thm:peps_fourier_regional_ranges} gives each local
    range inclusion. Every output slice is a sum over exterior input
    configurations of these block images, so all parent constraints
    are preserved.
\end{proof}

\begin{theorem}[Canonical parents in different representation coordinates]%
    \label{thm:peps_fourier_coordinate_parent}
    \lean{TNLean.PEPS.graphNumberedBondMatrix_coordinates_rightInverse,
        TNLean.PEPS.graphNumberedBondMatrix_coordinates_injective,
        TNLean.PEPS.map_regionParentGroundSpace_coordinates_eq,
        TNLean.PEPS.canonicalCoordinateParentEquiv,
        TNLean.PEPS.finrank_regionParentGroundSpace_coordinates_eq,
        TNLean.PEPS.map_ker_regionParentHamiltonian_coordinates_eq}
    \leanok
    Suppose $Q^\dagger Q=I$ and $L(g)=QU(g)Q^\dagger$.
    For supplied incident enumerations and any family of regions,
    \begin{align}
        \mathcal Q\mathcal P_{\mathcal R}(U) & =\mathcal P_{\mathcal R}(L),
                                             & \mathcal Q:\mathcal P_{\mathcal R}(U) & \xrightarrow{\ \cong\ }
        \mathcal P_{\mathcal R}(L).
        \label{eq:peps_fourier_coordinate_parent}
    \end{align}
    Their dimensions agree. The global operation associated to
    $Q^\dagger$ is its inverse on the entire physical space:
    $QQ^\dagger=I$ suffices for the right-inverse identity, and
    $Q^\dagger Q=I$ suffices for injectivity.
    For a finite region family, arbitrary positive local interactions
    whose kernels are the respective regional ranges satisfy
    $\mathcal Q(\ker H)=\ker K$.
    No vertex or edge coverage assumption is needed.
\end{theorem}
\begin{proof}\leanok
    \uses{thm:peps_fourier_exterior_blocks,
        thm:peps_fourier_reverse_coordinates,
        thm:peps_fourier_product_composition,thm:peps_region_parent_common_kernel}
    Regroup the numbered physical coordinates by edges.
    Composition of the global products reduces to
    $B(Q)B(Q^\dagger)=B(QQ^\dagger)=I$ at each edge.
    Reversing the order proves injectivity when $Q^\dagger Q=I$.
    Theorem~\ref{thm:peps_fourier_reverse_coordinates} supplies both
    cancellations and the representation identity for $Q^\dagger$.
    Apply Theorem~\ref{thm:peps_fourier_exterior_blocks} in both
    directions and restrict the inverse operations to the parent spaces.
    Positivity identifies the Hamiltonian kernels with those spaces.
\end{proof}

\begin{definition}[Dimension restoration followed by Fourier transformation]%
    \label{def:peps_fourier_semiregular_parent_map}
    \lean{TNLean.PEPS.semiRegularFourierParentMap,
        TNLean.PEPS.canonicalSemiRegularFourierParentEquiv}
    \leanok
    \uses{thm:peps_copy_parent_equivalence,thm:peps_fourier_coordinate_parent}
    Let $d_i>0$, and suppose an isometric matrix
    $Q:\C^{\coprod_i\{0,\ldots,d_i-1\}^2}\to\C^G$ satisfies
    $L(g)=QD^{[d]}(g)Q^\dagger$.
    Given incident enumerations at all three stages, define
    $T=\mathcal Q M_d$. When every edge is internal to some parent
    region, its restriction is the composition of the copy-restoration
    and Fourier parent-space isomorphisms.
\end{definition}

\begin{theorem}[The weighted canonical parent in the regular group basis]%
    \label{thm:peps_fourier_semiregular_parent}
    \lean{TNLean.PEPS.map_regionParentGroundSpace_semiRegular_fourier_eq,
        TNLean.PEPS.finrank_regionParentGroundSpace_semiRegular_fourier_eq,
        TNLean.PEPS.map_ker_regionParentHamiltonian_semiRegular_fourier_eq}
    \leanok
    \uses{def:peps_fourier_semiregular_parent_map}
    With the hypotheses of
    Definition~\ref{def:peps_fourier_semiregular_parent_map}, let $A_d$
    be the averaging tensor for $D=\bigoplus_iD_i$ with weights
    $d_i^{1/4}$, and let $A_{\mathrm{reg}}$ be the unweighted averaging
    tensor for $L$ in group coordinates. Then
    \begin{align}
        T\mathcal P_{\mathcal R}(A_d) & =
        \mathcal P_{\mathcal R}(A_{\mathrm{reg}}),
                                      & \dim\mathcal P_{\mathcal R}(A_d) & =
        \dim\mathcal P_{\mathcal R}(A_{\mathrm{reg}}).
    \end{align}
    For a finite region family and arbitrary positive local interactions
    with these exact regional kernels, $T(\ker H)=\ker K$.
\end{theorem}
\begin{proof}\leanok
    \uses{thm:peps_copy_parent_equivalence,thm:peps_fourier_coordinate_parent,
        thm:peps_region_parent_common_kernel}
    Copy restoration identifies the full weighted single-copy parent
    space with the unweighted $D^{[d]}$ parent space. The Fourier
    isomorphism then identifies that entire space with the regular
    parent space. Compose the two maps, use preservation of dimension,
    and identify each positive Hamiltonian kernel with its parent space.
\end{proof}

% This supplied-Fourier comparison is distinct from deriving the blocks and
% Fourier matrix of an arbitrary semi-regular U. That derivation and the
% physical G-injective comparison appear in the following section.
